\documentclass[11pt]{article}

\usepackage[margin=1in]{geometry}
\usepackage{graphicx}
\usepackage{booktabs}
\usepackage{amsmath}
\usepackage{amssymb}
\usepackage{hyperref}
\usepackage{caption}
\usepackage{subcaption}
\usepackage{xcolor}
\usepackage{placeins}
\usepackage{multirow}

\title{KABC Benchmark Results: Experiment 1 (Gaussian Mixtures),\\ Experiment 2 (Identical Means, Different Shapes), and\\ Experiment 3 (MovieLens-1M)}
\author{Kernel-based Active Clustering of Bandit Arms\\
\small Thuot, Vogt, Ghoshdastidar, Verzelen (arXiv:2601.07535)}
\date{\today}

\begin{document}
\maketitle

\noindent This report summarizes the experimental setup, baseline algorithms,
results, and conclusions for Experiments 1, 2, and 3 of the \texttt{kernel-active-bandit}
benchmark suite. Full design rationale lives in \texttt{docs/experiments.md}; this
document summarizes what was actually run and what the real (non-smoke-test)
results show.

\paragraph{Run provenance.} \texttt{results/exp1/results.parquet} holds 10{,}400 rows
(13 grid points $\times$ 2 budget multipliers $\times$ 4 algorithms $\times$ 100 reps,
$\approx$24.4 hours wall-clock). \texttt{results/exp2/results.parquet} holds 10{,}500
rows (21 grid points across three instance families, $\approx$19 hours wall-clock).
\texttt{results/exp3/results.parquet} holds 550 rows
(2 grid points, 100 reps per algorithm except TaS-FW's 50, $\approx$66 minutes
wall-clock). All three were generated by \texttt{scripts/run\_experiment.py} using
calibrated threshold constants (Section~\ref{sec:calibration}). Figures were
regenerated from these files by \texttt{notebooks/01\_gaussian\_mixtures.ipynb},
\texttt{notebooks/02\_same\_mean\_shapes.ipynb}, and
\texttt{notebooks/03\_movielens.ipynb}.

\tableofcontents

\section{Experimental Setup}

\subsection{Experiment 1 --- Gaussian mixtures: competitiveness in the parametric regime}

\paragraph{Claim tested.} In the mean-separated Gaussian setting (home turf of
mean-based methods like ACB), KABC's budget is within a modest factor of ACB's;
budget scales as $s_*^{-2}\log(N/\delta)$ (the paper's Theorem 3.1).

\paragraph{Instance family.} \texttt{GaussianMixtureEnv}: $N$ arms in $K=3$ balanced
clusters, arm $i$ in cluster $k$ distributed as $\mathcal{N}(\mu_k, \sigma^2 I_d)$,
$\sigma=1$. Cluster means $\mu_k = \frac{\Delta}{\sqrt2}(e_{2k-1}+e_{2k})$, giving each
pair of cluster means an actual Euclidean separation of $\Delta\sqrt2$.

\paragraph{Grid.} One knob swept at a time, others held at the base config
$N{=}60, K{=}3, d{=}10, \Delta{=}3$:
\begin{itemize}
  \item $\Delta \in \{1, 1.5, 2, 3, 4, 6\}$ (6 points --- the headline sweep)
  \item $d \in \{6, 10, 50, 100\}$ (4 points; $d=2$ excluded, construction requires $d \ge 2K=6$)
  \item $N \in \{30, 60, 120\}$ (3 points)
\end{itemize}
13 distinct grid points total (the base config is each sweep's own midpoint, so it
is run identically under all three sweep tags --- a documented, accepted
redundancy).

\paragraph{Budget regimes.} Every grid point above is run under two multipliers:
\begin{itemize}
  \item $\text{budget\_multiplier}=1.0$: the $\delta$-PAC \emph{guaranteed budget},
    $\text{budget\_cap} = 200\, N\, s_*^{-2}\log(N/\delta)$, $s_*$ computed numerically
    per grid point.
  \item $\text{budget\_multiplier}=0.25$: a new \emph{restricted budget} pass, same
    grid, budget cut to a quarter. Motivation: real deployments of active
    clustering often cannot afford the full guaranteed-budget amount; this
    measures graceful degradation under a hard $4\times$ cut, complementary to
    (not a replacement for) the guaranteed-budget story.
\end{itemize}
This gives $13 \times 2 = 26$ total grid points.

\paragraph{Protocol.} $\delta=0.05$ throughout. 100 Monte Carlo repetitions per
(grid point, algorithm). $\text{seed} = 1000 \cdot \text{exp\_id} + \text{rep}$. Pilot
pulls ($n_{\text{pilot}}=5$/arm, median-heuristic bandwidth estimation) count toward
the reported $\tau$.

\paragraph{Algorithms.} \texttt{kabc}, \texttt{kabc-sg}, \texttt{acb}, \texttt{unif-kkm}
(Section~\ref{sec:baselines}). \texttt{tas-fw} is \emph{not} run here --- its
finite-alphabet assumption does not hold for continuous Gaussian instances.

\subsection{Experiment 2 --- Identical means, different shapes: nonparametric separation}
\label{sec:setup-exp2}

\paragraph{Claim tested.} Mean-based active clustering methods (ACB) cannot
separate clusters that share a mean but differ in higher-order distributional
structure (variance, shape, or geometry); KABC's kernel/MMD machinery can.
This is the synthetic complement to Experiment 3's real-data version of the
same claim.

\paragraph{Three instance families, $N=40$, $K=2$ throughout:}
\begin{itemize}
  \item \textbf{\texttt{scale\_pair}} (2a): cluster A $=\mathcal N(0,1)$, cluster B
    $=\mathcal N(0,\sigma_B^2)$ --- same mean, different \emph{variance}. Swept:
    $\sigma_B \in \{1.5, 2, 3, 4\}$.
  \item \textbf{\texttt{moment\_matched\_mixture}} (2b): cluster A $=\mathcal
    N(0,1)$, cluster B $= 0.5\,\mathcal N(-a,b^2) + 0.5\,\mathcal N(a,b^2)$,
    $b^2=1-a^2$ --- same mean \emph{and} variance, different \emph{shape}
    (unimodal $\to$ bimodal as $a\to1$). Swept: $a$, with an
    \textbf{asymmetric range by budget regime} (see below).
  \item \textbf{\texttt{ring\_blob}} (2c): cluster A $=$ a 2D ring (radius 2,
    Gaussian noise $\sigma=0.1$), cluster B $=\mathcal N(0,I_2)$, both centered
    at the origin --- same mean, different \emph{2D geometry}. No sweep (one
    configuration).
\end{itemize}

\paragraph{Budget regimes and the \texttt{moment\_matched\_mixture} asymmetry.}
As in Experiment 1, every family is run at both
$\text{budget\_multiplier}=1.0$ (guaranteed budget, $\text{budget\_cap} =
200\,N\,s_*^{-2}\log(N/\delta)$) and $0.25$ (restricted budget). For
\texttt{moment\_matched\_mixture} specifically, the two regimes cover
\emph{different} $a$-ranges: guaranteed budget covers only $a \in
\{0.8,0.85,0.9,0.95\}$, while restricted budget covers the full $a \in
\{0.5,0.6,0.7,0.8,0.85,0.9,0.95\}$. This is deliberate, not missing data: direct
diagnostic runs (this session, not part of the main grid) showed \texttt{kabc}
hits an engineering safeguard --- a Gram-matrix memory guard added to
\texttt{kabc.py}'s doubling-trick loop to bound worst-case memory use --- on
most seeds for $a\in\{0.5,0.7\}$ at \emph{either} budget regime (the guard
trips on Gram-matrix size at a given iteration, not on remaining budget, so
shrinking the budget does not avoid it). \texttt{kabc} never reaches a real
$\delta$-PAC stopping decision at these points, so a $\tau$/error comparison
there would only be reporting safeguard trips, not convergence behavior ---
excluded from the guaranteed-budget pass for that reason. They remain in the
restricted-budget pass because ARI on the best partition-so-far is still
well-defined even when \texttt{stopped=False} (\S\ref{sec:calibration}'s
sibling generalization in \texttt{docs/experiments.md} \S0.4), giving the only
view in this report of the full $a$-range's degradation.

\paragraph{Protocol.} $\delta=0.05$. 100 reps per (family, grid point,
algorithm). \texttt{kabc-oracle} (needs the true $s_*^2$ as an oracle input)
is included only where \texttt{budget\_cap} is $s_*$-formula-derived, i.e.\
every Experiment 2 grid point.

\paragraph{Algorithms.} \texttt{kabc}, \texttt{kabc-sg}, \texttt{kabc-oracle},
\texttt{acb}, \texttt{unif-kkm}. \texttt{acb} is expected to be mean-blind and
never stop on every family here (all three are same-mean by construction), by
design, not as a failure.

\paragraph{Calibration reuse.} As elsewhere, \texttt{kabc}/\texttt{kabc-sg}/\texttt{acb}
reuse the single frozen calibration from Section~\ref{sec:calibration}
(calibrated against Experiment 1's \texttt{gaussian\_mixture} family) rather
than being re-calibrated per family. Section~\ref{sec:results-exp2} shows this
matters: \texttt{kabc-sg}'s calibrated constant, in particular, generalizes
poorly to all three Experiment 2 families.

\subsection{Experiment 3 --- MovieLens: real rating distributions}

\paragraph{Claim tested.} The ``same mean, different shape'' phenomenon that
mean-based methods cannot see (demonstrated on synthetic data in Experiment 2)
also occurs in real data (polarizing vs.\ consensus movies), and KABC exploits it
where mean-based methods fail.

\paragraph{Data.} MovieLens-1M (\texttt{data/ml-1m/}), filtered to movies with
$\ge 1000$ ratings.

\paragraph{Prototype selection} (deterministic, scripted --- \texttt{MovieLensEnv}):
\begin{enumerate}
  \item For each retained movie, compute its empirical rating distribution on
    $\{1,\ldots,5\}$ and pairwise $\mathrm{MMD}^2$ (RBF kernel, median-heuristic
    bandwidth).
  \item Find the pair maximizing $\mathrm{MMD}^2$ subject to $|\text{mean gap}| \le
    0.05$ --- the ``polarizing vs.\ consensus'' same-mean pair.
  \item Add $K-2$ further prototypes greedily maximizing minimum pairwise
    $\mathrm{MMD}^2$.
  \item $K=4$ prototypes total, 10 arms/prototype ($N=40$, main instance). Pulling
    an arm draws one rating with replacement from its prototype's pool.
\end{enumerate}

On the real data, the same-mean pair selected was \emph{Clear and Present Danger}
(1994, mean rating 3.73) vs.\ \emph{South Park: Bigger, Longer and Uncut} (1999,
mean rating 3.76) --- a 3-star-consensus family drama against a bimodal
love-it-or-hate-it comedy: essentially the same mean rating, very different
distribution shape (Figure~\ref{fig:exp3-fig1}, left panel).

A separate, much smaller \emph{reduced} instance ($K=3$ prototypes, 2 arms each,
$N=6$) is built so TaS-FW's finite-alphabet assumption holds natively (MovieLens
ratings \emph{are} a finite alphabet, $\{1,\ldots,5\}$), letting it run as a fair
baseline.

\paragraph{Budget.} Main instance: fixed $\text{budget\_cap}=2{,}000{,}000$ pulls (no
$s_*$-based formula --- $s_*$ is not in closed form for real, empirical rating
distributions). Reduced instance: $\text{budget\_cap}=20{,}000$.

\paragraph{Protocol.} $\delta=0.05$. 100 reps/algorithm (50 for TaS-FW). No
budget-multiplier sweep for Experiment 3.

\paragraph{Algorithms.} Main instance --- \texttt{kabc}, \texttt{kabc-sg},
\texttt{acb}, \texttt{unif-kkm}. Reduced instance --- \texttt{tas-fw}, \texttt{kabc}.

\subsection{Experiment 4 --- MovieLens random pairs: robustness across the real-data pair distribution}

\paragraph{Claim tested.} Experiment 3 uses one fixed, cherry-picked $K=4$
MovieLens instance. Experiment 4 asks whether KABC's advantage (and
ACB's/TaS-FW's relative behavior) is an artifact of that one instance, or
holds up when the movie pair is randomly resampled every repetition across
the real diversity of the MovieLens catalog, under three pairing regimes
that vary how much mean separation is present and whether the second movie
is cherry-picked (max $\mathrm{MMD}^2$) or sampled uniformly.

\paragraph{Data.} \texttt{data/ml-1m-207/} --- the 207 MovieLens-1M movies
with $\ge 1000$ ratings, pre-filtered so per-rep pair sampling does not
re-load/re-filter the full ratings file every repetition. Same discrete
rating support $\{1,\ldots,5\}$ and RBF-on-ratings kernel as Experiment 3,
with the same fixed median-heuristic bandwidth (bandwidth $=1.0$) computed
once over the full 207-movie pool.

\paragraph{Instance.} $N=8$ arms, $K=2$ clusters, 4 arms per prototype (movie).
For each repetition, movie $A$ is drawn uniformly at random from the 207
movies; movie $B$ is chosen from $A$ according to one of three
version-specific rules below. \texttt{MovieLensPairEnv} then resamples with
replacement from each movie's real rating pool to build that repetition's
arms. If no eligible $B$ exists for a drawn $A$, a fresh $A$ is redrawn
deterministically until a valid pair is found.

\paragraph{Version-specific pairing rules.}
\begin{description}
  \item[\texttt{exp4\_1} --- tight mean gap, MMD-maximizing.] Candidates: all
    movies with $|\text{mean}_A - \text{mean}_B| \le 0.1$. $B = \arg\max
    \mathrm{MMD}^2$ among candidates --- the same ``polarizing vs.\ consensus,
    cherry-picked for maximum shape difference at matched mean'' construction
    as Experiment 3, redrawn fresh every repetition.
  \item[\texttt{exp4\_2} --- moderate mean gap, uniform.] Candidates: all
    movies with $0.1 \le |\text{mean}_A - \text{mean}_B| \le 0.5$. $B$ sampled
    uniformly at random among candidates.
  \item[\texttt{exp4\_3} --- tight-ish mean gap, uniform.] Candidates: all
    movies with $|\text{mean}_A - \text{mean}_B| \le 0.2$. $B$ sampled
    uniformly at random among candidates --- an ablation of \texttt{exp4\_1}
    that keeps a similarly tight mean gap but removes the MMD-maximizing
    cherry-pick.
\end{description}

\paragraph{Budget.} Fixed $\text{budget\_cap}=20{,}000$ pulls for every
algorithm (not the $s_*$-formula type used elsewhere --- not meaningful at
this instance scale). $\delta=0.05$, matching every other experiment.
Calibrated $c_{\text{cal}}$ constants (\texttt{kabc}, \texttt{kabc-sg},
\texttt{acb}) are reused from the existing frozen calibration, per the
``calibrate once, reuse everywhere'' protocol --- Experiment 4's instance
family (random MovieLens pairs) is distributionally different from what
those constants were calibrated against, so any resulting miscalibration is
reported as a finding, not silently re-tuned.

\paragraph{Algorithms.} \texttt{kabc}, \texttt{kabc-sg}, \texttt{acb},
\texttt{tas-fw}, and two non-adaptive uniform-sampling baselines,
\texttt{unif-kkm-1x}/\texttt{unif-kkm-2x}, whose \texttt{budget\_cap} is set
\emph{per repetition} to $\min(\tau_{\texttt{kabc}}, 20{,}000)$ and
$\min(2\tau_{\texttt{kabc}}, 20{,}000)$ respectively, using that same
repetition's \texttt{kabc} run --- the sharpest possible comparison of
active vs.\ non-adaptive sampling at matched (or doubled) budget. 40
repetitions per algorithm per version (\texttt{tas-fw} included at full
scale; $N=8$ was chosen specifically to keep \texttt{tas-fw} tractable, see
\S\ref{sec:conclusions}).

\paragraph{A note on graceful degradation.} During this run, a real failure
mode was found and fixed: KABC's median-heuristic bandwidth estimation can
receive a pilot sample whose pairwise distances are all zero (e.g.\ a
5-pulls-per-arm pilot draw from a heavily-skewed movie's rating pool, by
chance containing only tied high-rating values). \texttt{median\_heuristic\_bandwidth}
now falls back to the \emph{mean} pairwise distance in this case (still
informative provided fewer than 100\% of pairs are tied), only raising if
the mean is also exactly zero, in which case KABC stops gracefully
(\texttt{stopped=False}) rather than crashing. This is the same convention
already used elsewhere in \texttt{kabc.py} for budget-too-tight and
memory-guard graceful stops.

\FloatBarrier
\section{Baseline Algorithms}
\label{sec:baselines}

\subsection{KABC and KABC-SG (this paper's own methods)}

KABC (Algorithm~2 of \cite{kabc2026}) runs a doubling-trick outer loop: at
iteration $k$, pull $n_k = \lceil 2^k \log(8(N^2-N)/\delta_k)\rceil$ fresh samples
per arm ($\delta_k = \delta/(4k^2)$), then build a same-cluster graph over all arm
pairs via a pairwise MMD-based test --- an edge $\{i,j\}$ is added iff the
empirical MMD statistic between arms $i$ and $j$ falls under a threshold. The
recovered partition is the graph's connected components; the loop stops the
first time the number of connected components equals the true $K$.

Two threshold variants are compared (identical outer loop, different pairwise
test):
\begin{description}
  \item[\texttt{kabc}] variance-aware, Wolfer--Alquier-style threshold (Eq.~6 of
    \cite{kabc2026}): sizes each pair's own threshold from that pair's
    \emph{empirical} RKHS variance, so a noisier pair automatically gets more
    slack.
  \item[\texttt{kabc-sg}] kernel-uniform sub-Gaussian threshold (the Discussion
    section's ablation in \cite{kabc2026}): uses only the fixed kernel bound
    (RBF kernel sup $=1$), identical for every pair regardless of that pair's
    own spread --- the ablation whose practical fragility Section~\ref{sec:conclusions}
    discusses in detail.
\end{description}

\subsection{ACB --- Active Clustering with Bandit feedback}

Source: \cite{acb2024}. ACB* (this repo's adaptive variant, ``$\theta_*$ and
$\Delta_*$ unknown,'' referred to throughout as ACB) is a \emph{mean}-based active
clustering procedure (SRI + ADC subroutines): it estimates each arm's mean
directly (no kernel/MMD machinery) and clusters by Euclidean distance between
empirical means, adaptively refining which arms to pull based on estimated
separation. It is the ``home turf'' competitor for Experiment~1 (mean-separated
Gaussians, exactly ACB's design setting) and the mean-blindness foil for
Experiment~3 (same-mean, different-shape data, where a mean-only method has no
signal by construction).

\subsection{TaS-FW --- Track-and-Stop with Frank-Wolfe}

Source: \cite{tasfw2024}. Track-and-Stop based on Frank-Wolfe optimization, for
$N$-ary clustering of arms whose distributions live on a common \emph{finite}
alphabet. Only run on Experiment~3's reduced ($N=6$) instance, where MovieLens
ratings on $\{1,\ldots,5\}$ satisfy the finite-alphabet assumption natively; not
run on Experiment~1's continuous Gaussian instances, which would require an
off-label discretization approximation outside the algorithm's proven
guarantees.

\subsection{Unif-KKM --- uniform sampling + kernel $k$-means (folklore baseline)}

No paper --- a standard, non-adaptive baseline implemented directly per the
design doc. Pulls a fixed number of samples per arm (sized to the same
budget cap the adaptive methods are given), then runs ordinary kernel
$k$-means on the resulting empirical kernel mean embeddings. Has no
confidence-based stopping rule and no $\delta$-PAC guarantee --- it always
spends its entire allotted budget and always ``completes'' ($\text{stopped}=
\text{True}$ by construction), which is why it is included purely as a
non-adaptive reference point, not as a $\delta$-correct competitor.

\FloatBarrier
\section{Calibration}
\label{sec:calibration}

\texttt{kabc}, \texttt{kabc-sg}, and \texttt{acb} each have a theoretically-derived
but conservative leading constant in their stopping threshold. Per the
calibration protocol (\texttt{docs/experiments.md} \S0.2), each constant was
empirically calibrated once (binary search for the smallest constant keeping
empirical error $\le \delta/2$), then frozen and reused unchanged across every
experiment. The frozen values actually used for this report:

\begin{table}[h]
\centering
\begin{tabular}{lccc}
\toprule
Algorithm & $c_{\text{cal}}$ & Calibration-set error & Held-out error \\
\midrule
\texttt{kabc} & 0.103516 & 0.000 & 0.000 \\
\texttt{kabc-sg} & 0.121094 & 0.000 & 0.000 \\
\texttt{acb} & 0.658984 & 0.020 & 0.010 \\
\texttt{tas-fw} & --- (no calibratable constant) & --- & --- \\
\bottomrule
\end{tabular}
\end{table}

\paragraph{Important caveat.} Calibration instances are a \emph{single} fixed
Gaussian-mixture instance (\texttt{GaussianMixtureEnv}'s construction is
deterministic given $N,K,d,\Delta$ --- there is no way to draw ``20 different
instances,'' only 20 independently-seeded noise draws of the one instance). The
frozen constants above were validated against that one instance's own real
budget regime, \emph{not} against Experiment~1's full, more varied grid.
Section~\ref{sec:conclusions} shows this matters directly: \texttt{kabc-sg}'s
and \texttt{acb}'s calibration does not fully generalize to every real
Experiment~1 grid point. \texttt{kabc-sg}'s frozen value (0.121094) was itself a
deliberate, documented override of the raw binary-search answer (0.117578),
chosen for margin above a demonstrated instability rather than pure
search-set efficiency.

\FloatBarrier
\section{Results --- Experiment 1}

\subsection{Guaranteed budget ($\text{budget\_multiplier}=1.0$): does the theory hold?}

Aggregated over all 13 grid points, 100 reps each (1{,}300 runs/algorithm):

\begin{table}[h]
\centering
\begin{tabular}{lccc}
\toprule
Algorithm & Stopped rate & Error rate & Mean ARI \\
\midrule
\texttt{kabc}     & 100.0\% & 0.0\%  & 1.000 \\
\texttt{unif-kkm} & 100.0\% & 0.0\%  & 1.000 \\
\texttt{acb}      & 100.0\% & 10.1\% & 0.987 \\
\texttt{kabc-sg}  &  65.5\% & 34.5\% & 0.870 \\
\bottomrule
\end{tabular}
\end{table}

\begin{figure}[h]
\centering
\includegraphics[width=0.7\textwidth]{../results/figures/exp1_fig1_tau_vs_delta.png}
\caption{Headline figure: $\log\tau$ vs.\ $\log\Delta$, all four methods, with the
theoretical $\tau \propto s_*^{-2}$ slope overlaid (dotted).}
\label{fig:exp1-fig1}
\end{figure}

\begin{figure}[h]
\centering
\begin{subfigure}{0.48\textwidth}
\includegraphics[width=\textwidth]{../results/figures/exp1_fig2_tau_vs_d.png}
\caption{$\tau$ vs.\ ambient dimension $d$.}
\end{subfigure}
\hfill
\begin{subfigure}{0.48\textwidth}
\includegraphics[width=\textwidth]{../results/figures/exp1_fig2b_tau_vs_N.png}
\caption{$\tau$ vs.\ number of arms $N$.}
\end{subfigure}
\caption{Budget scaling along the two secondary sweep dimensions.}
\label{fig:exp1-fig2}
\end{figure}

\begin{figure}[h]
\centering
\includegraphics[width=0.55\textwidth]{../results/figures/exp1_fig3b_error_heatmap.png}
\caption{Error rate at every grid point $\times$ algorithm $\times$ budget
regime. Asterisks mark cells exceeding $\delta=0.05$.}
\label{fig:exp1-fig3b}
\end{figure}

Per-grid-point detail (guaranteed budget), the specific points where error
exceeds $\delta=0.05$:
\begin{itemize}
  \item \texttt{kabc-sg} exceeds $\delta$ at: $\Delta{=}1$ (97\%), $\Delta{=}1.5$
    (84\%), $\Delta{=}2$ (45\%), $N{=}30$ (16\%), $d{=}50$ (100\%), $d{=}100$ (100\%).
  \item \texttt{acb} exceeds $\delta$ at: $\Delta{=}1$ (12\%), $\Delta{=}1.5$
    (11\%), $\Delta{=}2$ (13\%), $N{=}120$ (7\%), $d{=}50$ (45\%), $d{=}100$ (39\%).
  \item \texttt{kabc} and \texttt{unif-kkm} exceed $\delta$: nowhere.
\end{itemize}

\subsection{Restricted budget ($\text{budget\_multiplier}=0.25$): graceful degradation?}

\begin{figure}[h]
\centering
\includegraphics[width=0.6\textwidth]{../results/figures/exp1_fig5_restricted_budget_ari_vs_delta.png}
\caption{Mean ARI at $0.25\times$ budget vs.\ $\Delta$.}
\label{fig:exp1-fig5}
\end{figure}

\begin{table}[h]
\centering
\caption{Mean ARI at $0.25\times$ budget, by $\Delta$.}
\begin{tabular}{lcccccc}
\toprule
$\Delta$ & 1.0 & 1.5 & 2.0 & 3.0 & 4.0 & 6.0 \\
\midrule
\texttt{kabc}     & 1.000 & 1.000 & 1.000 & 1.000 & 1.0 & 1.0 \\
\texttt{kabc-sg}  & 0.880 & 0.917 & 0.960 & 0.996 & 1.0 & 1.0 \\
\texttt{acb}      & 0.991 & 0.989 & 0.987 & 0.999 & 1.0 & 1.0 \\
\texttt{unif-kkm} & 1.000 & 1.000 & 1.000 & 1.000 & 1.0 & 1.0 \\
\bottomrule
\end{tabular}
\end{table}

\begin{figure}[h]
\centering
\begin{subfigure}{0.48\textwidth}
\includegraphics[width=\textwidth]{../results/figures/exp1_fig5b_restricted_budget_ari_vs_d.png}
\caption{Mean ARI at $0.25\times$ budget vs.\ $d$.}
\end{subfigure}
\hfill
\begin{subfigure}{0.48\textwidth}
\includegraphics[width=\textwidth]{../results/figures/exp1_fig5c_restricted_budget_ari_vs_N.png}
\caption{Mean ARI at $0.25\times$ budget vs.\ $N$.}
\end{subfigure}
\caption{Restricted-budget degradation along the two secondary sweep dimensions.}
\label{fig:exp1-fig5bc}
\end{figure}

\begin{table}[h]
\centering
\caption{Mean ARI at $0.25\times$ budget, by $d$ and by $N$.}
\begin{tabular}{lcccc}
\toprule
$d$ & 6 & 10 & 50 & 100 \\
\midrule
\texttt{kabc}     & 1.000 & 1.000 & 1.000 & 1.000 \\
\texttt{kabc-sg}  & 1.000 & 0.996 & 0.385 & 0.026 \\
\texttt{acb}      & 1.000 & 0.999 & 0.944 & 0.930 \\
\texttt{unif-kkm} & 1.000 & 1.000 & 1.000 & 1.000 \\
\bottomrule
\end{tabular}
\quad
\begin{tabular}{lccc}
\toprule
$N$ & 30 & 60 & 120 \\
\midrule
\texttt{kabc}     & 1.000 & 1.000 & 1.000 \\
\texttt{kabc-sg}  & 0.943 & 0.996 & 1.000 \\
\texttt{acb}      & 0.999 & 0.999 & 0.998 \\
\texttt{unif-kkm} & 1.000 & 1.000 & 1.000 \\
\bottomrule
\end{tabular}
\end{table}

\FloatBarrier
\section{Results --- Experiment 2}
\label{sec:results-exp2}

\begin{figure}[h]
\centering
\includegraphics[width=0.95\textwidth]{../results/figures/exp2_fig0_qualitative_overview.png}
\caption{What ``identical means, different shape'' looks like in each of the
three families: same variance-mismatch densities (2a), same mean/variance
shape-mismatch densities colored by $a$ (2b), and the ring/blob 2D geometry
(2c). Purely illustrative reconstructions of each family's instance
geometry, not algorithm output.}
\label{fig:exp2-fig0}
\end{figure}

\subsection{\texttt{moment\_matched\_mixture} (2b)}

\begin{figure}[h]
\centering
\includegraphics[width=0.85\textwidth]{../results/figures/exp2_fig1_moment_matched_headline.png}
\caption{Guaranteed budget: same-mean/variance densities at $a=0.8$ (left);
stopping budget $\tau$ vs.\ $a$ with ACB's ARI on the right axis (right).
\texttt{kabc-oracle} never converges at any grid point (\S\ref{sec:results-exp2}
below).}
\label{fig:exp2-fig1}
\end{figure}

\begin{figure}[h]
\centering
\includegraphics[width=0.65\textwidth]{../results/figures/exp2_fig1b_restricted_budget_ari_vs_a.png}
\caption{Restricted budget ($0.25\times$), full $a$ range including the
$a\in\{0.5,0.6,0.7\}$ points excluded from the guaranteed-budget pass.}
\label{fig:exp2-fig1b}
\end{figure}

At guaranteed budget ($a\in\{0.8,\ldots,0.95\}$), \texttt{kabc} converges on
$98$--$100\%$ of reps ($\tau$ falling from $\approx\!90$k to $\approx\!39$k
pulls as $a$ grows), \texttt{kabc-sg} converges on $94$--$100\%$, and
\texttt{acb}'s ARI stays flat at $0.00$ throughout (mean-blind, as expected).
\texttt{kabc-oracle} converges \textbf{0\% of the time at every single grid
point in this experiment} --- its one-shot required budget
($n_{\text{pilot}}+N\cdot n^*$, see \texttt{kabc.py}'s \texttt{fit\_oracle}
docstring) exceeds \texttt{budget\_cap} everywhere under this experiment's
$\text{constant}=200$ formula. This is a real finding about the oracle
baseline's budget requirement at this $N/K/\delta$, not a bug or a
smoke-test artifact (\S\ref{sec:conclusions}).

The restricted-budget view (Figure~\ref{fig:exp2-fig1b}) shows a sharp phase
transition: at $a\in\{0.5,0.6\}$ \emph{every} algorithm scores ARI
$\approx 0$ --- the shape gap is too small to detect with only a quarter of
the guaranteed budget, regardless of adaptivity. At $a=0.8$, \texttt{kabc}
jumps immediately to ARI $=1.00$ while \texttt{unif-kkm} (naive, non-adaptive)
catches up more gradually ($0.60\to 1.00$ by $a=0.9$), and \texttt{kabc-sg}
lags well behind both ($0.19\to 0.47$) even at $a=0.95$. \texttt{acb} stays at
ARI $\approx 0$ across the entire range.

\subsection{\texttt{scale\_pair} (2a)}

\begin{figure}[h]
\centering
\includegraphics[width=0.6\textwidth]{../results/figures/exp2_fig2_scale_pair.png}
\caption{Guaranteed budget: stopping budget $\tau$ vs.\ $\sigma_B$, ACB's ARI
inset.}
\label{fig:exp2-fig2}
\end{figure}

\begin{figure}[h]
\centering
\includegraphics[width=0.55\textwidth]{../results/figures/exp2_fig2b_restricted_budget_ari_vs_sigma.png}
\caption{Restricted budget ($0.25\times$) ARI vs.\ $\sigma_B$.}
\label{fig:exp2-fig2b}
\end{figure}

\texttt{scale\_pair} is markedly easier than \texttt{moment\_matched\_mixture}:
\texttt{kabc} converges $100\%$ of the time at every $\sigma_B$ (guaranteed
budget), and both \texttt{kabc} and \emph{even naive} \texttt{unif-kkm} reach
ARI $=1.00$ at every $\sigma_B$ under the restricted budget too ---
variance mismatches are detectable with far less budget than the
higher-order shape mismatches of 2b. \texttt{kabc-sg} converges only
$63$--$93\%$ of the time (guaranteed budget) and its restricted-budget ARI
ranges $0.44$--$0.92$. \texttt{acb} and \texttt{kabc-oracle} both converge
$0\%$ of the time, as expected for a same-mean family / an infeasible
one-shot budget respectively.

\subsection{\texttt{ring\_blob} (2c)}

\begin{figure}[h]
\centering
\includegraphics[width=0.85\textwidth]{../results/figures/exp2_fig4_ring_blob.png}
\caption{Left: raw samples ($N=40$, true cluster), 5 samples/arm. Right:
\texttt{kabc}'s recovered partition, 1 (noisy) sample/arm for a readable
scatter --- the ARI is annotated directly since a single-sample plot does not
visually reconstruct a clean ring shape even at perfect recovery.}
\label{fig:exp2-fig4}
\end{figure}

\begin{table}[h]
\centering
\caption{\texttt{ring\_blob}: mean ARI, standard deviation of ARI, mean
$\tau$ (budget spent), and fraction of reps that reached a genuine
$\delta$-PAC stop, for every algorithm, at both budget regimes ($n=100$
reps/row).}
\begin{tabular}{lccccc}
\toprule
Algorithm & Budget regime & Mean ARI & Std ARI & Mean $\tau$ & Stopped rate \\
\midrule
\texttt{kabc}        & guaranteed (1.0$\times$)  & 1.000 & 0.000 &  18{,}815 & 100\% \\
\texttt{kabc-sg}     & guaranteed (1.0$\times$)  & 0.989 & 0.101 &  17{,}385 &  99\% \\
\texttt{kabc-oracle} & guaranteed (1.0$\times$)  & 0.000 & 0.000 &     200 &   0\% \\
\texttt{acb}         & guaranteed (1.0$\times$)  & 0.000 & 0.000 & 543{,}732 &   0\% \\
\texttt{unif-kkm}    & guaranteed (1.0$\times$)  & 1.000 & 0.000 & 543{,}760 & 100\% \\
\midrule
\texttt{kabc}        & restricted (0.25$\times$) & 1.000 & 0.000 &  18{,}815 & 100\% \\
\texttt{kabc-sg}     & restricted (0.25$\times$) & 0.987 & 0.101 &  15{,}572 &  98\% \\
\texttt{kabc-oracle} & restricted (0.25$\times$) & 0.000 & 0.000 &     200 &   0\% \\
\texttt{acb}         & restricted (0.25$\times$) & 0.000 & 0.000 & 131{,}663 &   0\% \\
\texttt{unif-kkm}    & restricted (0.25$\times$) & 1.000 & 0.000 & 131{,}680 & 100\% \\
\bottomrule
\end{tabular}
\label{tab:exp2-ringblob}
\end{table}

\texttt{kabc} recovers the true ring-vs-blob partition perfectly and
identically under both budget regimes (\texttt{kabc}'s own $\tau\approx19$k
is already well under the restricted $0.25\times$ cap, so cutting the budget
has no effect on it). \texttt{unif-kkm} also reaches perfect ARI at both
regimes, but only by spending its \emph{entire} budget every time (mean
$\tau\approx$ the budget cap itself) --- this 2D geometric separation is easy
enough that even non-adaptive uniform sampling with enough total pulls
recovers it exactly, unlike the harder \texttt{moment\_matched\_mixture}
cases above. \texttt{acb} and \texttt{kabc-oracle} both fail completely (ARI
$=0$), for the two structurally different reasons already established: \texttt{acb} is
mean-blind on this same-mean family; \texttt{kabc-oracle} never affords its
required one-shot budget.

\subsection{Price of adaptivity}

\begin{figure}[h]
\centering
\includegraphics[width=0.45\textwidth]{../results/figures/exp2_fig3_price_of_adaptivity.png}
\caption{$\tau_{\texttt{kabc}}/\tau_{\texttt{kabc-oracle}}$ across all
Experiment 2 configurations.}
\label{fig:exp2-fig3}
\end{figure}

This figure is included for completeness but is uninformative in this report:
since \texttt{kabc-oracle} converges $0\%$ of the time at every grid point
(above), \emph{zero} of the $2{,}100$ (grid point, algorithm, rep) pairs have
both \texttt{kabc} and \texttt{kabc-oracle} converged, so no reliable
price-of-adaptivity ratio is computable from this run. Populating this figure
would require re-running Experiment 2 with a larger \texttt{budget\_cap}
constant than $200$ specifically for \texttt{kabc-oracle} --- at the cost of
no longer sharing \texttt{budget\_cap} with \texttt{kabc}/\texttt{kabc-sg}/\texttt{acb} on
the same grid point, and thus not included here.

\FloatBarrier
\section{Results --- Experiment 3}

\begin{figure}[h]
\centering
\includegraphics[width=0.85\textwidth]{../results/figures/exp3_fig1_same_mean_pair_and_ari.png}
\caption{Left: rating histograms of the real same-mean movie pair. Right: final
ARI per algorithm, annotated with mean $\tau$.}
\label{fig:exp3-fig1}
\end{figure}

\begin{table}[h]
\centering
\caption{Main instance ($N=40$, $K=4$ prototypes), 100 reps/algorithm.}
\begin{tabular}{lcccc}
\toprule
Algorithm & Stopped rate & Mean $\tau$ & Error rate & Mean ARI \\
\midrule
\texttt{kabc}     & 100.0\% & 19{,}672     & 0.0\%   & 1.000 \\
\texttt{kabc-sg}  &  71.0\% & 69{,}326     & 63.0\%  & 0.861 \\
\texttt{acb}      &   0.0\% & 1{,}999{,}901 & 100.0\% & 0.000 \\
\texttt{unif-kkm} & 100.0\% & 2{,}000{,}000 & 0.0\%   & 1.000 \\
\bottomrule
\end{tabular}
\end{table}

\begin{table}[h]
\centering
\caption{Reduced instance ($N=6$): TaS-FW vs.\ KABC.}
\begin{tabular}{lccccc}
\toprule
Algorithm & Stopped rate & Mean $\tau$ & Error rate & Mean ARI & Mean wall-clock/rep \\
\midrule
\texttt{tas-fw} & 100.0\% & 4{,}357 & 0.0\% & 1.000 & 39.64\,s \\
\texttt{kabc}   &  98.0\% & 1{,}852 & 2.0\% & 0.995 & 0.23\,s \\
\bottomrule
\end{tabular}
\end{table}

\FloatBarrier
\section{Results --- Experiment 4}

\begin{figure}[h]
\centering
\includegraphics[width=\textwidth]{../results/figures/exp4_fig1_ari_bars.png}
\caption{Mean ARI ($\pm$ std) per algorithm, 3 pairing versions $\times$ 6 algorithms, 40 reps each.}
\label{fig:exp4-fig1}
\end{figure}

\begin{figure}[h]
\centering
\includegraphics[width=\textwidth]{../results/figures/exp4_fig2_tau_bars.png}
\caption{Mean $\tau$ ($\pm$ std, log scale) per algorithm.}
\label{fig:exp4-fig2}
\end{figure}

\begin{figure}[h]
\centering
\includegraphics[width=\textwidth]{../results/figures/exp4_fig3_wallclock_bars.png}
\caption{Mean wall-clock/rep in seconds ($\pm$ std, log scale) per algorithm.}
\label{fig:exp4-fig3}
\end{figure}

\begin{table}[h]
\centering
\caption{Experiment 4, $N=8$, 40 reps/algorithm/version. Stopped rate, mean $\tau$, mean ARI.}
\begin{tabular}{llccc}
\toprule
Version & Algorithm & Stopped rate & Mean $\tau$ & Mean ARI \\
\midrule
\multirow{6}{*}{\texttt{exp4\_1}}
 & \texttt{kabc}        &  97.5\% &  6{,}749 & 0.942 \\
 & \texttt{kabc-sg}     &  65.0\% &  6{,}425 & 0.464 \\
 & \texttt{acb}         &   0.0\% & 19{,}994 & 0.000 \\
 & \texttt{tas-fw}      &  30.0\% & 18{,}448 & 1.000 \\
 & \texttt{unif-kkm-1x} & 100.0\% &  6{,}749 & 0.949 \\
 & \texttt{unif-kkm-2x} & 100.0\% & 11{,}964 & 0.949 \\
\midrule
\multirow{6}{*}{\texttt{exp4\_2}}
 & \texttt{kabc}        &  90.0\% &  7{,}296 & 0.937 \\
 & \texttt{kabc-sg}     &  55.0\% &  8{,}425 & 0.433 \\
 & \texttt{acb}         &  17.5\% & 19{,}153 & 0.089 \\
 & \texttt{tas-fw}      &  35.0\% & 18{,}166 & 1.000 \\
 & \texttt{unif-kkm-1x} & 100.0\% &  7{,}296 & 0.927 \\
 & \texttt{unif-kkm-2x} & 100.0\% & 12{,}294 & 0.881 \\
\midrule
\multirow{6}{*}{\texttt{exp4\_3}}
 & \texttt{kabc}        &  52.5\% & 12{,}127 & 0.488 \\
 & \texttt{kabc-sg}     &  87.5\% &  3{,}007 & 0.160 \\
 & \texttt{acb}         &   0.0\% & 19{,}993 & 0.000 \\
 & \texttt{tas-fw}      &   2.5\% & 19{,}920 & 0.987 \\
 & \texttt{unif-kkm-1x} & 100.0\% & 12{,}127 & 0.550 \\
 & \texttt{unif-kkm-2x} & 100.0\% & 18{,}123 & 0.574 \\
\bottomrule
\end{tabular}
\label{tab:exp4-summary}
\end{table}

\FloatBarrier
\section{Inferences and Conclusions}
\label{sec:conclusions}

\paragraph{6.1 KABC's $\delta$-PAC guarantee holds empirically, everywhere it was tested.}
Across all 26 Experiment~1 grid points (both budget regimes) and both
Experiment~3 instances, \texttt{kabc}'s empirical error was $0.0\%$ except a
single $2\%$ figure on Experiment~3's tiny $N=6$ reduced instance (still
comfortably under $\delta=0.05$). No other adaptive method in this benchmark
achieved that consistency. This is the paper's central claim, holding up
under a real, non-smoke-scale run across a genuinely varied grid and real
MovieLens data, not just one convenient configuration.

\paragraph{6.2 KABC's budget scales with separation roughly as the theory predicts.}
Figure~\ref{fig:exp1-fig1}'s dashed $s_*^{-2}$ reference line and \texttt{kabc}'s
actual curve track each other reasonably closely across the full $\Delta$
sweep, supporting Theorem~3.1's scaling claim empirically.

\paragraph{6.3 The variance-aware threshold is materially more robust than the
kernel-uniform sub-Gaussian ablation --- direct empirical evidence for the
paper's own claimed benefit.}
\texttt{kabc-sg} is the clear failure mode of this benchmark: $34.5\%$ aggregate
error under the guaranteed budget (vs.\ \texttt{kabc}'s $0.0\%$), and a
catastrophic collapse to near-random clustering (mean ARI $0.026$) at $d=100$
under a restricted budget. The mechanism, established by direct inspection
during calibration: \texttt{kabc-sg}'s threshold uses one identical cutoff for
every arm pair, so whichever single pair is noisiest for a given draw
determines success or failure for the \emph{whole} clustering; \texttt{kabc}'s
variance-aware threshold sizes each pair's own cutoff from that pair's
empirical spread, absorbing exactly this kind of pair-specific noise. Every
\texttt{kabc-sg} failure inspected was a near-miss (one or two ``straggler''
arms failing to join their true cluster, ARI $0.86$--$0.98$ typically) that ran
out of budget rather than resolving --- consistent with a uniformly
slightly-too-tight threshold rather than a qualitative break.

\paragraph{6.4 ACB, on its own home turf, does not reliably satisfy its own
confidence level at harder configurations --- a calibration-generalization
gap, not a design flaw.}
\texttt{acb}'s aggregate error ($10.1\%$) exceeds $\delta=0.05$, concentrated at
small $\Delta$ ($11$--$13\%$) and large $d$ ($39$--$45\%$). Unlike \texttt{kabc-sg},
\texttt{acb}'s stopped rate is $100\%$ everywhere --- its errors are confident,
completed wrong answers, not budget-exhaustion near-misses. Given \texttt{acb}'s
calibration-set/held-out error ($2.0\%/1.0\%$) looked fine in isolation, this is
best read as the single fixed calibration instance failing to surface the
harder Experiment~1 grid points, not as evidence against ACB's design. This
should \emph{not} be conflated with the mean-blindness story below, a
qualitatively different and more severe failure mode.

\paragraph{6.5 On real data, ACB is not just imprecise --- it is completely blind
to same-mean, different-shape structure, exactly as the theory predicts.}
Experiment~3's headline number is unambiguous: \texttt{acb}'s mean ARI on the
real same-mean movie pair is \emph{exactly} $0.000$, with $0\%$ of runs ever
stopping (it exhausts its full $2{,}000{,}000$-pull budget every repetition).
This is qualitatively different from the calibration gap in 6.4: a mean-based
method has no signal whatsoever when two clusters share a mean, by
construction, regardless of budget or calibration. \texttt{kabc} recovers the
partition perfectly (ARI $=1.000$) using under $1\%$ of \texttt{acb}'s budget
($19{,}672$ vs.\ $1{,}999{,}901{+}$ pulls). This directly confirms, on real data,
the paper's core motivating claim: mean-based active clustering methods are
blind to shape-only cluster differences that occur in practice.

\paragraph{6.6 KABC-SG's fragility from \S6.3 is not an Experiment-1-specific
artifact --- it reproduces on three structurally distinct synthetic families.}
A direct calibration-generalization check this session (holding the frozen
\texttt{kabc-sg} constant fixed, testing empirical error against 100 fresh
reps on each Experiment~2 family) found error rates of $35\%$
(\texttt{scale\_pair}), $66\%$ (\texttt{moment\_matched\_mixture}), and $4\%$
(\texttt{ring\_blob}) --- all far above the required $\delta/2=2.5\%$ that
held on the calibration instance itself, while \texttt{kabc}'s constant
passed cleanly ($0\%$ error) on all three. This is exactly the mechanism from
\S6.3: \texttt{kabc-sg}'s single, non-adaptive cutoff (fixed $\bar g=1$) has
no way to tighten itself against a specific family's variance/shape
structure the way \texttt{kabc}'s per-pair empirical threshold does, so it
transfers poorly to families qualitatively different from the Gaussian
instance it was calibrated on. This is independent evidence for the same
underlying claim, not a restatement of it.

\paragraph{6.7 KABC-oracle's required one-shot budget is infeasible under
Experiment 2's budget formula --- a real, reportable finding about the
oracle baseline, not a run failure.} At every single Experiment~2 grid point
(all 21), \texttt{kabc-oracle} converges $0\%$ of the time: its required
one-shot pull count ($n_{\text{pilot}} + N n^*$) exceeds
$\text{budget\_cap}$ under the shared $\text{constant}=200$ formula used by
every algorithm on a given grid point. Because \texttt{kabc} and
\texttt{kabc-oracle} therefore never both converge on the same rep, the
``price of adaptivity'' figure (Figure~\ref{fig:exp2-fig3}) is uninformative
in this report --- not a plotting bug, but a direct consequence of the
oracle needing a substantially larger budget than an adaptive method ever
needs to spend, at this $N/K/\delta$.

\paragraph{6.8 Nonparametric separability has a sharp difficulty gradient
across shape vs.\ variance vs.\ geometry, and a genuinely hard region where
budget alone does not help.} \texttt{scale\_pair}'s variance mismatch is
separable by every adaptive method (and even naive \texttt{unif-kkm}) at
full ARI even under a $4\times$ budget cut. \texttt{moment\_matched\_mixture}
is markedly harder: at $a\le0.6$, \emph{every} algorithm --- including
\texttt{kabc} and generously-budgeted \texttt{unif-kkm} --- scores ARI
$\approx0$ under the restricted budget; this is a genuinely hard instance
region, not an adaptivity gap, since more budget without more separation
does not help a non-adaptive method either. Only past $a=0.8$ does
\texttt{kabc}'s advantage become visible: it reaches ARI $=1.00$ immediately
where \texttt{unif-kkm} needs until $a\approx0.9$ to catch up, and
\texttt{kabc-sg} never fully catches up within this $a$ range at all.
\texttt{ring\_blob}'s 2D geometric mismatch sits at the easy end again ---
even \texttt{unif-kkm} reaches perfect ARI, provided it is given its entire
(large) budget.

\paragraph{6.9 Restricted-budget degradation is graceful for KABC (absent, in
fact) and ACB, but sharp and dimension-driven for KABC-SG.}
Under a $4\times$ budget cut, \texttt{kabc}'s mean ARI stays at exactly $1.000$
across every one of the 13 grid points --- no measurable degradation. \texttt{acb}
degrades mildly and smoothly ($0.93$--$0.99$ in the worst cases). \texttt{kabc-sg}
is the outlier: fine on the $\Delta$ and $N$ sweeps ($\ge 0.88$ and $\ge 0.94$
respectively) but collapsing specifically as ambient dimension grows ($0.385$
at $d=50$, $0.026$ at $d=100$ --- indistinguishable from random). This
dimension-specific collapse (not mirrored in the other sweeps) suggests
\texttt{kabc-sg}'s fixed, non-adaptive threshold interacts badly with
higher-dimensional Gram-matrix statistics specifically, compounding the
pair-uniformity fragility of \S6.3.

\paragraph{6.10 TaS-FW is a fair, correctness-competitive baseline at the one
scale its assumptions hold, at a real, quantified compute cost.}
On Experiment~3's $N=6$ reduced instance, both \texttt{tas-fw} and \texttt{kabc}
achieve essentially perfect recovery (ARI $1.000$ vs.\ $0.995$) with comparable,
modest budgets ($4{,}357$ vs.\ $1{,}852$ pulls). The real divergence is wall-clock
cost per repetition: \texttt{tas-fw} averages $39.64$\,s, \texttt{kabc} averages
$0.23$\,s --- roughly $172\times$ more compute per decision, a real, measured
confirmation of the design doc's ``orders of magnitude more compute''
characterization.

\paragraph{6.12 ACB's mean-blindness on same-mean data (\S6.5) is confirmed
again, now averaged over a random sample of the real pair distribution
rather than one fixed instance.} \texttt{acb}'s mean ARI is \emph{exactly}
$0.000$ in both \texttt{exp4\_1} and \texttt{exp4\_3} (tight mean-gap
regimes, $0\%$ ever stopping), and only $0.089$ in \texttt{exp4\_2}
(moderate gap, $17.5\%$ stopped) --- consistent with \S6.5's single-instance
finding, now shown to hold across a random sample of movie pairs rather than
one cherry-picked example. Every other method (\texttt{kabc},
\texttt{tas-fw}, both \texttt{unif-kkm} variants) reaches mean ARI $\ge
0.88$ on the same \texttt{exp4\_1}/\texttt{exp4\_2} pairs, so this is not an
artifact of instance difficulty.

\paragraph{6.13 KABC-SG's fragility (\S6.3, \S6.6) reproduces a third time,
now on real discrete rating data, and \texttt{exp4\_3} shows KABC itself is
not immune to hard pairs.} \texttt{kabc-sg}'s mean ARI never exceeds $0.464$
on any Experiment~4 version, well below \texttt{kabc}'s $0.942$/$0.937$ on
the matching \texttt{exp4\_1}/\texttt{exp4\_2} versions --- the same
single-non-adaptive-threshold mechanism from \S6.3. \texttt{exp4\_3} (tight
mean gap, no MMD-maximizing cherry-pick) is markedly harder for
\emph{every} adaptive method, including \texttt{kabc} itself (mean ARI
drops to $0.488$, stopped rate $52.5\%$): with the cherry-pick removed, many
sampled same-mean pairs have genuinely small shape difference, and
$\text{budget\_cap}=20{,}000$ is not always enough to resolve them. This
shows \texttt{exp4\_1}'s strong \texttt{kabc} result partly depends on the
MMD-maximizing selection rule, not only on the mean gap being tight.

\paragraph{6.14 TaS-FW remains correctness-competitive but at an even
starker, now precisely quantified, compute cost (extends \S6.10).} At
$N=8$, \texttt{tas-fw} reaches mean ARI $\ge 0.987$ on every version
(highest or tied-highest of all six algorithms throughout), but at a mean
wall-clock cost of $789$--$908$\,s/rep --- two to three orders of magnitude
slower than every other algorithm (all under $2.1$\,s/rep). Unlike every
other algorithm, \texttt{tas-fw} rarely reports $\text{stopped}=\text{True}$
(2.5\%--35\% across versions) yet still lands its partition estimate
correctly well before its formal stopping rule fires, given the shared
$20{,}000$-pull budget --- its accuracy is not in question, only its
practicality at scale (\S\ref{sec:conclusions}, ``Why $N$ cannot grow''
consideration carried over from \texttt{docs/experiments.md}).

\paragraph{6.11 Open caveats for the write-up.}
\begin{itemize}
  \item \texttt{kabc-sg}'s and \texttt{acb}'s calibrated constants do not fully
    generalize from the single fixed calibration instance to Experiment~1's
    full real grid (\S\ref{sec:calibration}, \S6.4), and \texttt{kabc-sg}'s
    constant additionally fails to generalize to all three Experiment~2
    families (\S6.6). Any claim about their error rates should be scoped to
    ``as calibrated in this repo,'' not read as an indictment of the
    underlying algorithms' theoretical guarantees.
  \item Experiment~2's ``price of adaptivity'' figure
    (Figure~\ref{fig:exp2-fig3}) is uninformative as run, since
    \texttt{kabc-oracle} never affords its required one-shot budget at any
    grid point (\S6.7); populating it would require a separate, larger
    \texttt{budget\_cap} just for \texttt{kabc-oracle}, not attempted here.
  \item Figure~4 of the original design (appendix: theoretical vs.\ calibrated
    constants panel, Experiment 1) requires a second, uncalibrated ($c=1$)
    comparison run that has not been generated; not included here.
  \item Experiment~4's fixed $\text{budget\_cap}=20{,}000$ interacts with
    KABC's doubling-trick schedule in a way that is worth flagging: at
    $N=8$, $\delta=0.05$, the schedule's cumulative cost after 6 rounds is
    $13{,}832$, and the 7th round alone would cost a further $14{,}728$
    (total $28{,}560$) --- comfortably past the $20{,}000$ cap. Every
    Experiment~4 \texttt{kabc}/\texttt{kabc-sg} repetition that failed to
    converge stopped at exactly this $13{,}832$ boundary (or the version's
    analogous one), one full doubling round short of what it would need to
    attempt the next round at all --- a budget-cap artifact of this specific
    schedule/cap combination, not a property of the algorithm's difficulty
    on those pairs. A future run with a larger cap (chosen to actually admit
    at least one more full doubling round) would give KABC/KABC-SG a fairer
    chance to convergence on the harder pairs.
\end{itemize}

\begin{thebibliography}{9}

\bibitem{kabc2026}
Thuot, Vogt, Ghoshdastidar, Verzelen.
\textit{Kernel-based Active Clustering of Bandit Arms}.
arXiv:2601.07535, 2026.

\bibitem{acb2024}
Thuot, Carpentier, Giraud, Verzelen.
\textit{Active clustering with bandit feedback}.
arXiv:2406.11485, 2024.

\bibitem{tasfw2024}
Yavas, Huang, Tan, Scarlett.
\textit{A General Framework for Clustering and Distribution Matching with Bandit Feedback}.
arXiv:2409.05072, 2024.

\end{thebibliography}

\end{document}
